\chapter{Conclusiones y Trabajo Futuro}

\section{Conclusiones}

Este trabajo abordó la automatización de la traducción de escenarios BDD en lenguaje natural hacia representaciones IBDD, una etapa crítica del pipeline BDD-a-MBT que, en la práctica, suele requerir intervención manual experta. Para ello se diseñó e implementó un enfoque híbrido que combina un módulo de traducción basado en LLM con validación sintáctica determinística mediante gramática formal IBDD y parser Earley, complementado por un ciclo iterativo de corrección sobre casos fallidos.

Desde el punto de vista de la ingeniería, el trabajo aporta: (i) un pipeline reproducible, con proveedor/modelo, prompt y dataset parametrizados; (ii) una gramática operacional de IBDD integrada a un validador automático; (iii) un mecanismo de corrección asistido por diagnóstico de errores sintácticos; y (iv) una implementación organizada en etapas desacopladas, con resultados intermedios persistidos en archivos JSON, lo que favorece trazabilidad, reproducibilidad y ejecución independiente de cada componente.

En términos empíricos, la evaluación final sobre el conjunto \textit{test} de 90 escenarios mostró que el sistema alcanza niveles altos de validez sintáctica final en las cuatro configuraciones analizadas, con medias entre 84.22\,\% y 91.11\,\%. Además, la variabilidad entre corridas se mantuvo acotada para la métrica principal, lo que respalda la estabilidad del comportamiento observado bajo el protocolo experimental adoptado.

Respecto de las preguntas de investigación:

\begin{description}
  \item[RQ1:] El sistema logra traducir automáticamente escenarios BDD a IBDD con alta tasa de conformidad sintáctica en el conjunto \textit{test}, medida de forma objetiva mediante el parser.
  \item[RQ2:] El idioma del prompt muestra una asociación más marcada con la validez inicial que el idioma del dataset; luego del ciclo de corrección, las diferencias entre configuraciones se reducen.
  \item[RQ3:] El ciclo iterativo de corrección aporta mejoras consistentes en todas las configuraciones (ganancias positivas de SVR y recuperación de casos inicialmente fallidos), aunque en promedio utiliza el máximo de rondas configurado.
\end{description}

En síntesis, los resultados sostienen que la combinación de generación con LLM y validación/corrección formal constituye una estrategia viable para automatizar la etapa BDD$\rightarrow$IBDD en términos de validez sintáctica, dentro del contexto experimental evaluado. Al mismo tiempo, la evidencia obtenida delimita con claridad el alcance actual del sistema: bajo el protocolo evaluado se observa alta validez sintáctica, mientras que la fidelidad semántica de la traducción permanece como dimensión abierta.

\section{Trabajo Futuro}

La primera línea de evolución propuesta es explorar una extensión del pipeline actual hacia un esquema con agentes especializados. En problemas donde la resolución requiere múltiples decisiones interdependientes, un flujo estrictamente lineal ofrece menor capacidad para representar estrategias de búsqueda, verificación cruzada y corrección iterativa. En este marco, agentes especializados podrían asumir funciones diferenciadas (traducción, validación, diagnóstico, propuesta de corrección y control de calidad), coordinados por una política de orquestación explícita. Como extensión concreta, esta arquitectura debería formalizarse y documentarse de manera sistemática, especificando roles, contratos de entrada y salida, criterios de activación entre agentes y registro de decisiones, con el objetivo de preservar reproducibilidad y auditabilidad. Bajo este esquema, el sistema podría incorporar ciclos de verificación cruzada entre agentes, mecanismos de consenso para resolver conflictos y estrategias de reintento guiadas por diagnósticos intermedios.

El sistema desarrollado en este trabajo implementa validación sintáctica de las traducciones IBDD mediante una gramática formal y un parser, pero no incorpora una etapa de validación semántica automática que verifique si la traducción preserva correctamente el significado del escenario BDD original. Esto constituye una limitación de alcance: una traducción puede ser sintácticamente correcta y al mismo tiempo semánticamente deficiente; por ejemplo, podría omitir una entidad mencionada en el escenario, utilizar un gate de salida donde el BDD describe una acción del usuario, o perder un valor numérico relevante para una postcondición.

Durante el desarrollo de este trabajo se consideraron alternativas preliminares para abordar parcialmente este problema, por ejemplo mediante heurísticas léxicas (cotejo de entidades, acciones y valores numéricos entre el texto BDD y la traducción IBDD), verificaciones de consistencia numérica, uso del LLM como apoyo para revisión semántica, o validadores apoyados en reglas de dominio. Sin embargo, no se avanzó con una implementación de estas estrategias, ya que su diseño y evaluación exceden el alcance de este trabajo. En particular, la naturaleza abreviada de IBDD, donde variables como \texttt{JF} representan entidades como ``job file'' y gates como \texttt{?submit} condensan frases completas, dificulta el cotejo directo con el lenguaje natural. Por ello, la validación semántica permanece como una línea de trabajo futuro abierta.

Otra dirección de trabajo futuro es la generalización de la capa de integración con proveedores LLM. Si bien la arquitectura actual opera con proveedor y modelo parametrizados por configuración y desacopla el resto del pipeline del cliente concreto, la implementación disponible soporta explícitamente dos proveedores: OpenAI (utilizado en las corridas principales tras la selección del modelo en el Capítulo~\ref{sec:seleccion-modelo}) y Ollama (incluido para habilitar ejecución local de modelos de código abierto en contextos donde se prioriza soberanía de datos). Extender el sistema a otros proveedores (por ejemplo, Anthropic) requiere implementar un adaptador adicional en la clase \texttt{LLMClient}, en particular para resolver diferencias en autenticación, formato de mensajes, manejo de errores y soporte de salidas estructuradas. Una evolución natural del sistema consiste en formalizar esta capa como un conjunto de adaptadores intercambiables con capacidades declaradas, de modo de aproximar una interoperabilidad más completa entre proveedores sin modificar los módulos de traducción, análisis de errores y orquestación.

Una línea de trabajo futuro con impacto directo es la evaluación del pipeline completo BDD$\rightarrow$IBDD$\rightarrow$BDDTS$\rightarrow$STS$\rightarrow$Tests. En este trabajo, la evaluación se centró en la validez sintáctica de la traducción a IBDD; como extensión, resulta relevante medir cómo esos resultados se propagan a las etapas formales posteriores y a la calidad de los casos de prueba generados. Esto permitiría incorporar métricas de extremo a extremo, por ejemplo proporción de escenarios que llegan a tests ejecutables, cobertura alcanzada por los tests sobre el modelo STS y tasa de fallas atribuibles a errores introducidos en la etapa de traducción. Esta evaluación integral ofrecería evidencia más fuerte sobre la utilidad práctica del sistema en un flujo MBT completo.
